\documentclass[aspectratio=169,10pt]{beamer}
\usetheme[numbering=fraction,progressbar=frametitle]{metropolis}
\usepackage{fontspec}
\setsansfont{Latin Modern Sans}
\usepackage{amsmath,amssymb,mathtools,graphicx,booktabs}
\usepackage{polyglossia}
\setdefaultlanguage{english}
\setbeamertemplate{navigation symbols}{}
\setbeamerfont{frametitle}{size=\large}
\setbeamerfont{normal text}{size=\normalsize}
\setlength{\parskip}{.45em}
\hypersetup{pdfauthor={},pdfsubject={Lecture slides},pdfkeywords={},pdftitle={Conditional probability}}
\providecommand{\N}{\mathbb{N}}
\providecommand{\Z}{\mathbb{Z}}
\providecommand{\D}{\mathbb{D}}
\providecommand{\Q}{\mathbb{Q}}
\providecommand{\R}{\mathbb{R}}
\providecommand{\C}{\mathbb{C}}
\providecommand{\vect}[1]{\overrightarrow{#1}}
\providecommand{\vi}{\vec{\imath}}
\providecommand{\vj}{\vec{\jmath}}
\providecommand{\vk}{\vec{k}}
\providecommand{\norm}[1]{\left\lVert #1\right\rVert}
\providecommand{\vc}[2]{\begin{pmatrix}#1\\#2\end{pmatrix}}
\providecommand{\deter}[2]{\det\left(#1,#2\right)}
\providecommand{\abs}[1]{\left\lvert #1\right\rvert}
\providecommand{\dd}{\,\mathrm{d}}
\providecommand{\e}{\mathrm{e}}
\providecommand{\Card}{\operatorname{Card}}
\providecommand{\E}{\operatorname{E}}
\providecommand{\V}{\operatorname{V}}
\providecommand{\degre}{^{\circ}}
\providecommand{\un}[1]{\,\mathrm{#1}}
\providecommand{\sci}[2]{#1\times10^{#2}}
\providecommand{\pgcd}{\operatorname{PGCD}}
\providecommand{\ppcm}{\operatorname{PPCM}}
\title{Conditional probability}
\subtitle{Lecture and worked examples}
\date{}
\author{}
\begin{document}
\begin{frame}[plain]
\vfill
{\usebeamerfont{title}\inserttitle\par}
\medskip
{\usebeamerfont{subtitle}\insertsubtitle\par}
\medskip
{\usebeamercolor[fg]{progress bar}\rule{\textwidth}{.5pt}}
\vfill
\end{frame}
\begin{frame}{Reference course}
In a weighted tree, the probabilities leaving the same node sum to \(1\). We multiply the probabilities along a path: \(P(A\cap B)=P(A)P_A(B)\). We add the probabilities of mutually exclusive paths leading to the event we are looking for.
\end{frame}
\begin{frame}{Reference course}
Do not confuse \(P_A(B)\) and \(P_B(A)\). The condition comes after the word “given” and is written as a subscript in this notation.
\end{frame}
\begin{frame}{Reference course}
\(40\%\) of customers buy online. Among them, \(12\%\) return their order; among the others, \(5\%\) return it. Then \(P(R)=0.40\times0.12+0.60\times0.05=0.078\). The overall return rate is therefore \(7.8\%\).
\end{frame}
\begin{frame}{Conditioning}
\begin{itemize}
\item Define a conditional probability.
\item Move from counts to probabilities.
\item Interpret.

\end{itemize}
\end{frame}
\begin{frame}{Conditioning}
In a survey, \(180\) customers out of \(600\) use a premium card and \(72\) of these \(180\) customers buy every week. Calculate \(P(P)\), \(P(P\cap H)\) and \(P_P(H)\), then interpret.
\end{frame}
\begin{frame}{Conditioning}
For \(P(A)>0\), \(P_A(B)=\dfrac{P(A\cap B)}{P(A)}\). This is the probability of \(B\) within the population restricted to \(A\).
\end{frame}
\begin{frame}{Conditioning}
Take the data of the investigation and write the solution in order: definition of the objects, property used, exact or approximate calculation, then a sentence giving the answer. The model answer appears in the teacher edition.
\end{frame}
\begin{frame}{Conditioning}
\textbf{One-hour plan:} 7 min of fluency; 10 min of pooling and written record; 28 min of essential exercises; 5 min of correction; 3 min of exit ticket.

\textbf{Two-hour plan:} add 20 min of investigation, 20 min of formalisation, a second independent exercise and 10 min of comparative correction.

\textbf{Differentiation:} give struggling students a four-box sheet (data, model, calculation, conclusion). Ask fast students to produce a counterexample, a justification or a second method.
\end{frame}
\begin{frame}{Weighted trees}
\begin{itemize}
\item Construct and complete a tree.
\item Read a conditional probability.
\item Check the sums.

\end{itemize}
\end{frame}
\begin{frame}{Weighted trees}
\begin{columns}[T,onlytextwidth]
\begin{column}{.57\textwidth}
\(35\%\) of orders are express. Among them, \(4\%\) are late; among standard orders, \(9\%\) are late. Construct and complete the tree.
\end{column}
\begin{column}{.4\textwidth}
\centering
\includegraphics[width=\linewidth,height=4.5cm,keepaspectratio]{figures/stmg-06-2-en.pdf}
\end{column}
\end{columns}
\end{frame}
\begin{frame}{Weighted trees}
In a tree, the branches leaving a node sum to \(1\). The second-level probabilities are conditional.
\end{frame}
\begin{frame}{Weighted trees}
Take the data of the investigation and write the solution in order: definition of the objects, property used, exact or approximate calculation, then a sentence giving the answer. The model answer appears in the teacher edition.
\end{frame}
\begin{frame}{Weighted trees}
\textbf{One-hour plan:} 7 min of fluency; 10 min of pooling and written record; 28 min of essential exercises; 5 min of correction; 3 min of exit ticket.

\textbf{Two-hour plan:} add 20 min of investigation, 20 min of formalisation, a second independent exercise and 10 min of comparative correction.

\textbf{Differentiation:} give struggling students a four-box sheet (data, model, calculation, conclusion). Ask fast students to produce a counterexample, a justification or a second method.
\end{frame}
\begin{frame}{Product of probabilities along a path}
\begin{itemize}
\item Calculate an intersection.
\item Combine paths.
\item Translate an event.

\end{itemize}
\end{frame}
\begin{frame}{Product of probabilities along a path}
\begin{columns}[T,onlytextwidth]
\begin{column}{.57\textwidth}
Using the tree from session 40, calculate the probability of an order being express and late, then the total probability of a delay. Interpret both results.
\end{column}
\begin{column}{.4\textwidth}
\centering
\includegraphics[width=\linewidth,height=4.5cm,keepaspectratio]{figures/stmg-06-3-en.pdf}
\end{column}
\end{columns}
\end{frame}
\begin{frame}{Product of probabilities along a path}
We multiply along a path: \(P(A\cap B)=P(A)P_A(B)\). We add the mutually exclusive paths that realise the event.
\end{frame}
\begin{frame}{Product of probabilities along a path}
Take the data of the investigation and write the solution in order: definition of the objects, property used, exact or approximate calculation, then a sentence giving the answer. The model answer appears in the teacher edition.
\end{frame}
\begin{frame}{Product of probabilities along a path}
\textbf{One-hour plan:} 7 min of fluency; 10 min of pooling and written record; 28 min of essential exercises; 5 min of correction; 3 min of exit ticket.

\textbf{Two-hour plan:} add 20 min of investigation, 20 min of formalisation, a second independent exercise and 10 min of comparative correction.

\textbf{Differentiation:} give struggling students a four-box sheet (data, model, calculation, conclusion). Ask fast students to produce a counterexample, a justification or a second method.
\end{frame}
\begin{frame}{Independence}
\begin{itemize}
\item Test independence.
\item Compare probabilities.
\item Justify a conclusion.

\end{itemize}
\end{frame}
\begin{frame}{Independence}
In a population, \(P(A)=0.4\), \(P(B)=0.3\) and \(P(A\cap B)=0.12\). Test independence in two ways. Then change \(P(A\cap B)\) to \(0.15\) and start again.
\end{frame}
\begin{frame}{Independence}
\(A\) and \(B\) are independent if \(P(A\cap B)=P(A)P(B)\), equivalently if \(P_A(B)=P(B)\) when \(P(A)>0\).
\end{frame}
\begin{frame}{Independence}
Take the data of the investigation and write the solution in order: definition of the objects, property used, exact or approximate calculation, then a sentence giving the answer. The model answer appears in the teacher edition.
\end{frame}
\begin{frame}{Independence}
\textbf{One-hour plan:} 7 min of fluency; 10 min of pooling and written record; 28 min of essential exercises; 5 min of correction; 3 min of exit ticket.

\textbf{Two-hour plan:} add 20 min of investigation, 20 min of formalisation, a second independent exercise and 10 min of comparative correction.

\textbf{Differentiation:} give struggling students a four-box sheet (data, model, calculation, conclusion). Ask fast students to produce a counterexample, a justification or a second method.
\end{frame}
\begin{frame}{Partition and total probability}
\begin{itemize}
\item Recognise a partition.
\item Apply total probability.
\item Organise the calculation.

\end{itemize}
\end{frame}
\begin{frame}{Partition and total probability}
\begin{columns}[T,onlytextwidth]
\begin{column}{.57\textwidth}
A factory has three lines A, B, C producing \(20\%\), \(50\%\) and \(30\%\) of the items respectively. Their defect rates are \(1\%\), \(2\%\) and \(4\%\). Calculate the overall defect rate.
\end{column}
\begin{column}{.4\textwidth}
\centering
\includegraphics[width=\linewidth,height=4.5cm,keepaspectratio]{figures/stmg-06-4-en.pdf}
\end{column}
\end{columns}
\end{frame}
\begin{frame}{Partition and total probability}
If \((A_i)\) is a partition, \(P(B)=\sum_i P(A_i)P_{A_i}(B)\). Each term corresponds to a path leading to \(B\).
\end{frame}
\begin{frame}{Partition and total probability}
Take the data of the investigation and write the solution in order: definition of the objects, property used, exact or approximate calculation, then a sentence giving the answer. The model answer appears in the teacher edition.
\end{frame}
\begin{frame}{Partition and total probability}
\textbf{One-hour plan:} 7 min of fluency; 10 min of pooling and written record; 28 min of essential exercises; 5 min of correction; 3 min of exit ticket.

\textbf{Two-hour plan:} add 20 min of investigation, 20 min of formalisation, a second independent exercise and 10 min of comparative correction.

\textbf{Differentiation:} give struggling students a four-box sheet (data, model, calculation, conclusion). Ask fast students to produce a counterexample, a justification or a second method.
\end{frame}
\begin{frame}{Problems in context}
\begin{itemize}
\item Choose a tree or a table.
\item Solve a problem in context.
\item Interpret.

\end{itemize}
\end{frame}
\begin{frame}{Problems in context}
Turn a table of counts (customers/subscription/cancellation) into a tree, then answer a question in each direction: probability of cancelling given the subscription; probability of having the subscription given the cancellation.
\end{frame}
\begin{frame}{Problems in context}
Choosing between a tree and a table: the tree highlights the order of conditioning; the table suits counts and reversed readings well.
\end{frame}
\begin{frame}{Problems in context}
Take the data of the investigation and write the solution in order: definition of the objects, property used, exact or approximate calculation, then a sentence giving the answer. The model answer appears in the teacher edition.
\end{frame}
\begin{frame}{Problems in context}
\textbf{One-hour plan:} 7 min of fluency; 10 min of pooling and written record; 28 min of essential exercises; 5 min of correction; 3 min of exit ticket.

\textbf{Two-hour plan:} add 20 min of investigation, 20 min of formalisation, a second independent exercise and 10 min of comparative correction.

\textbf{Differentiation:} give struggling students a four-box sheet (data, model, calculation, conclusion). Ask fast students to produce a counterexample, a justification or a second method.
\end{frame}
\begin{frame}{Mixed practice}
\begin{itemize}
\item Calculate a reverse probability.
\item Chain total probability and conditioning.
\item Check the result.

\end{itemize}
\end{frame}
\begin{frame}{Mixed practice}
\begin{columns}[T,onlytextwidth]
\begin{column}{.57\textwidth}
In the situation of session 43, given that an item is defective, calculate the probability that it comes from line C. Explain why this result is not \(30\%\).
\end{column}
\begin{column}{.4\textwidth}
\centering
\includegraphics[width=\linewidth,height=4.5cm,keepaspectratio]{figures/stmg-06-5-en.pdf}
\end{column}
\end{columns}
\end{frame}
\begin{frame}{Mixed practice}
For a reverse probability, first calculate the total probability of the conditioning event, then use the definition of conditional probability.
\end{frame}
\begin{frame}{Mixed practice}
Take the data of the investigation and write the solution in order: definition of the objects, property used, exact or approximate calculation, then a sentence giving the answer. The model answer appears in the teacher edition.
\end{frame}
\begin{frame}{Mixed practice}
\textbf{One-hour plan:} 7 min of fluency; 10 min of pooling and written record; 28 min of essential exercises; 5 min of correction; 3 min of exit ticket.

\textbf{Two-hour plan:} add 20 min of investigation, 20 min of formalisation, a second independent exercise and 10 min of comparative correction.

\textbf{Differentiation:} give struggling students a four-box sheet (data, model, calculation, conclusion). Ask fast students to produce a counterexample, a justification or a second method.
\end{frame}
\begin{frame}{Review and assessment}
\begin{itemize}
\item Use all the methods.
\item Write up with the notation.
\item Remedy mistakes.

\end{itemize}
\end{frame}
\begin{frame}{Review and assessment}
Complete a four-station circuit: complete a tree; calculate an intersection; apply the total probability formula; test independence. Each group leaves a written justification for the next group.
\end{frame}
\begin{frame}{Review and assessment}
Always name the events, complete all the branches, write the formula before calculating and interpret the final probability.
\end{frame}
\begin{frame}{Review and assessment}
Take the data of the investigation and write the solution in order: definition of the objects, property used, exact or approximate calculation, then a sentence giving the answer. The model answer appears in the teacher edition.
\end{frame}
\begin{frame}{Review and assessment}
\textbf{One-hour plan:} 7 min of fluency; 10 min of pooling and written record; 28 min of essential exercises; 5 min of correction; 3 min of exit ticket.

\textbf{Two-hour plan:} add 20 min of investigation, 20 min of formalisation, a second independent exercise and 10 min of comparative correction.

\textbf{Differentiation:} give struggling students a four-box sheet (data, model, calculation, conclusion). Ask fast students to produce a counterexample, a justification or a second method.
\end{frame}
\end{document}
